\subsection{METODOLOGÍA}

La primera etapa de la investigación se centro en la búsqueda de los distintos modelos distinados
a la verificación de hablantes. 

\textcolor{red}{En cuanto a este aparto, estuve trajando en la recopilación de bibliografía para determinar qué 
modelo aplicado a la identificación de hablantes es apropiado emplear. Sobre todo, qué tecnica es necesaria 
aplicar para la utilización del modelo en formato ligero para su ejecución dentro del microcontrolador.
Recién se pudo detallar que el modelo a emplear sea el ECAPA-TDNN, tanto con su versión para PC como su versión 
ligera, ECAPA-TDNNLite para microcontrolador. Según la bibliografía revisada, la idea es emplear uno de los dos para 
el enrolamiento de nuevos usuarios y el otro para la verificación de hablantes. Dado a que según el paper de referencia,
de esta manera se consigue menor EER, es decir, mayor precisión.
}